\section{Lean implementation and the route to object-theory proofs}
\label{sec:formalization}

\subsection{What the supplied Lean files establish}

The accompanying sources were compiled with Lean~4.32.0, using only the
bundled \texttt{Std} library. The file \texttt{BellCore.lean} defines finite
sums, the two-variable Stirling recurrence, Bell numbers as row sums, and
factorials. It proves the Stirling boundary identities, positivity of Bell
numbers, the weighted recurrence
\[
 B_{n+1}=\sum_{k=0}^{n}(k+1)S(n,k),
\]
monotonicity, the upper bounds $B_{n+1}\le(n+1)B_n$ and $B_n\le n!$, and
the values through $B_{10}=115975$. The numerical identities use
definitional reduction. The file \texttt{Audit.lean} prints the trusted-axiom
dependencies of nine representative theorems. The computational identities
and the column-one identity use no axioms; the other printed proofs use
Lean's \texttt{propext} and \texttt{Quot.sound}. None uses
\texttt{Classical.choice} or \texttt{sorryAx}.

These facts describe the delivered formalization exactly. In particular,
the package does \emph{not} contain a Lean-checked proof of the saddle-point
theorem, nor a constructed \texttt{ConciseIPRA.Prf} derivation of the Bell
theory. Ordinary Lean theorems about natural numbers are metatheorems, even
when their proofs are constructive. Constructing derivations in a specified
first-order theory is a separate task. The mathematics of the preceding
sections supplies the human proof; this section specifies the remaining
formalization work and supplies executable definitions for its main targets.

The companion file \texttt{BellTargets.lean} defines total arithmetic
programs and propositions with names ending in \texttt{Target}. Such a definition records a
goal; it does not assert that goal. There are no added axioms or placeholder
proofs. This distinction also applies to the proposed mixed-radix Bell
program and the approximation programs below: their definitions compile,
while their mathematical correctness proofs remain to be translated into
Lean.

\subsection{The inspected ProveIt interfaces}

The repository was inspected at commit
\texttt{c0d0dc37f60aea10c0182ba4eece9eae8c7e13c6}.
The relevant namespace is \texttt{LeanProofs.ConciseIPRA}.
The source root is
\begin{quote}\small
\path{Logic/PrimitiveRecursiveArithmetic/ConciseIPRA/Lean/ConciseIPRA}.
\end{quote}
The audit JSON records the Git blob identifiers of every inspected module.
The upstream project itself was not rebuilt for this package.

The
\href{https://github.com/VladimirReshetnikov/ProveIt/blob/c0d0dc37f60aea10c0182ba4eece9eae8c7e13c6/Logic/PrimitiveRecursiveArithmetic/ConciseIPRA/Lean/ConciseIPRA/PR.lean}{\texttt{PR.lean} interface}
contains the following exact objects.
\begin{itemize}
\item \texttt{PR k}: primitive-recursive symbols of arity $k$, generated by
\texttt{zero}, \texttt{succ}, \texttt{proj}, \texttt{comp}, and
\texttt{prec}. In the last constructor the recursion variable is the first
argument, and the step receives the counter, the previous value, and the
parameters.
\item \texttt{PTm n}: terms with $n$ free variables; substitution is
\texttt{PTm.sub} and its generic interface is \texttt{Syn}.
\item \texttt{PRAAx}: the defining equations, the two successor axioms, and
induction for formulas satisfying \texttt{Fm.QF}.
\item \texttt{PRA.precZero}, \texttt{PRA.precSucc}, and
\texttt{PRA.qfInd}: reusable constructors for the corresponding
object-theory derivations.
\end{itemize}
The
\href{https://github.com/VladimirReshetnikov/ProveIt/blob/c0d0dc37f60aea10c0182ba4eece9eae8c7e13c6/Logic/PrimitiveRecursiveArithmetic/ConciseIPRA/Lean/ConciseIPRA/Deduction.lean}{\texttt{Deduction.lean} interface}
defines \texttt{Prf Ax} as intuitionistic natural deduction with equality,
including hypotheses, instantiated theory axioms, introduction and
elimination rules for the connectives and quantifiers, and simultaneous
equality substitution. Its derived operations include \texttt{Prf.subst},
\texttt{Prf.trans}, \texttt{Prf.eqCongr}, and \texttt{Prf.eqTrans}.

The
\href{https://github.com/VladimirReshetnikov/ProveIt/blob/c0d0dc37f60aea10c0182ba4eece9eae8c7e13c6/Logic/PrimitiveRecursiveArithmetic/ConciseIPRA/Lean/ConciseIPRA/BiInterpretation.lean}{\texttt{BiInterpretation.lean} interface}
provides \texttt{pra\_iff\_teq} and \texttt{pra\_theorem\_iff}; the
\href{https://github.com/VladimirReshetnikov/ProveIt/blob/c0d0dc37f60aea10c0182ba4eece9eae8c7e13c6/Logic/PrimitiveRecursiveArithmetic/ConciseIPRA/Lean/ConciseIPRA/NativeBiInterpretation.lean}{native interface}
provides \texttt{pra\_iff\_native}. These translate an already constructed
$\ipra$ derivation to the concise scalar or native calculus. They do not
turn a semantic Bell-number theorem into a derivation automatically. The
iterative implementation can also use \texttt{Tm.prc} and its proved
equations \texttt{Q.prc\_zero}, \texttt{Q.prc\_succ} from
\texttt{TeqRec.lean}.

\subsection{First-order codes and bounded invariants}

All finite objects may be encoded by natural numbers. For definiteness,
one may use
\[
 \langle a,b\rangle=2^a(2b+1)-1.
\]
The projections search only below the input code and are primitive
recursive. Encode a signed integer by a sign and a magnitude, a rational
by a signed numerator and a positive denominator, and a finite array by
its length together with a mixed-radix word. Addition, multiplication,
comparison, division with remainder, factorial, powers, bounded sums,
bounded products, and bounded searches then have fixed PR definitions.
Invalid inputs receive specified default values, so every function is total.
Theorems about valid codes carry validity hypotheses.

An especially direct first-order Bell program fixes the input $N$ and sets
\[
 A=(N+2)^{N+2}+1,\qquad w=N+1,\qquad
 d_k(c)=\left\lfloor c/A^k\right\rfloor\bmod A.
\]
Start with $c_0=1$, and perform $N$ steps
\[
 c_{j+1}=\sum_{k=0}^{N}
 \begin{cases}
 0,&k=0,\\
 k\,d_k(c_j)+d_{k-1}(c_j),&k>0
 \end{cases} A^k.
\]
Finally return $\sum_{k=0}^{N}d_k(c_N)$. These are the definitions
\texttt{digit}, \texttt{packedStep}, \texttt{packedRows}, and
\texttt{packedBell} in the delivered source.

The invariant at stage $j\le N$ says that the digits are $S(j,k)$ for
$0\le k\le N$, and that every digit is below $A$. There is no carry because
\[
 S(j,k)\le B_j\le j!\le (N+1)^{N+1}<A.
\]
The invariant is a bounded conjunction of elementary comparisons and
equalities. Replace that conjunction by its primitive-recursive Boolean
test and apply induction to the single equation saying that the test is
$1$. When the final unbounded symbol $S(n,k)$ is defined from a table whose
size is determined by $n+k$, first prove that enlarging the table width or
the base does not change earlier entries. This is a second bounded
simultaneous invariant. It yields the usual Stirling recursion with no
appeal to higher-type recursion.

This construction matters because the elegant Lean definition
\texttt{stirling : Nat $\to$ Nat $\to$ Nat} recurses over a row viewed as a
function. Its ordinary Lean proof of vanishing above the diagonal also
generalizes over the second argument during induction. Neither feature
should be silently identified with quantifier-free induction in the
object theory. The packed-table invariant is the explicit replacement.

For the combinatorial interpretation, a restricted-growth word of length
$n$ is encoded in base $n+1$ below $(n+1)^n$. Its validity test checks the
initial label and the bounded inequalities against preceding maximum
labels. The number of blocks is the maximum label plus one, with the empty
word treated separately. Counting valid words, and counting those with
exactly $k$ labels, are therefore bounded sums. Every bijection used in the
finite proofs becomes an explicitly bounded encoder and decoder together
with their inverse equations. Finite groups in the congruence argument are
handled in exactly this fashion: their actions are permutations of bounded
code sets, and orbit calculations are bounded arithmetic.

There is no uniform primitive-recursive evaluator for \emph{all} PR
programs assumed here. Each displayed algorithm is compiled to one fixed
PR definition; subsequent formulas contain that symbol. Generic lemmas
about a function $f$ in Lean are metatheoretic schemas, instantiated with
these fixed programs before constructing an object-theory proof.

\subsection{Rational approximation programs and their error estimates}
\label{subsec:lean-approximation}

The record \texttt{QCode} consists of an integer numerator and a natural
denominator. Its order test is cross multiplication. The operations do not
require normalization, and positivity of a denominator is a separate
arithmetic property. Write $\varepsilon_t=2^{-t}$.

\paragraph{Exponential.}
For a valid rational code $x=u/v$, put $m=|u|+1$ and
\[
 K=8(m+1)^2+2t+16,\qquad
 E_t(x)=\sum_{j=0}^{K-1}\frac{x^j}{j!}.
\]
This is \texttt{expApprox}. It satisfies
\[
 |E_t(x)-e^x|\le 2^{-t}.
\]
Here and below the displayed real notation abbreviates the corresponding
finite Cauchy-tail inequalities. To prove the exponential bound, note that
$|x|\le m$ and, for $k\ge K$, the last half of the factors in $k!$ are at
least $\lfloor k/2\rfloor+1$. For $k=2h$ this gives
$m^k/k!\le(m^2/(h+1))^h\le 2^{-k}$; for $k=2h+1$ it gives the bound
$2^{1-k}$. Thus every finite tail starting at $K$ is bounded by
$\sum_{k\ge K}2^{1-k}=2^{2-K}\le2^{-t}$.
For nonnegative $x$, all summands are nonnegative, so $E_t(x)\le e^x$.
This one-sided property is used in the bounded search for $W$.

\paragraph{Logarithm.}
For $x=u/v>0$, with positive integers $u,v$, put
\[
 z=\frac{u-v}{u+v},\qquad
 M=(u+v)^2 2^{t+4},\qquad
 L_t(x)=2\sum_{j=0}^{M-1}\frac{z^{2j+1}}{2j+1}.
\]
This is \texttt{logApprox}. Its finite-tail bound is
\[
 \left|2\sum_{j=M}^{J}\frac{z^{2j+1}}{2j+1}\right|
 \le\frac{2}{(2M+1)(1-z^2)}
 =\frac{(u+v)^2}{2uv(2M+1)}
 \le2^{-t}.
\]
The logarithm identities used elsewhere follow by the ordinary finite
power-series manipulations with this tail bound, or by rational
integration of $1/x$. Invalid and nonpositive inputs are sent to zero in
the total program; no theorem interprets those outputs as logarithms.

\paragraph{Lambert $W$.}
Let $n\ge1$, $\delta=2^{-(t+2)}$, and test the dyadic grid
$x_k=k\delta$ for $0\le k\le(n+1)/\delta$. Use
\[
 s=t+n+8,
 \qquad x_kE_s(x_k)\ge n
\]
as the success test. Choose its first successful index by the explicit
bounded search \texttt{boundedMin}. The last test succeeds because
$E_s(n+1)\ge1$. Write $r=W(n)$ and $r_t$ for the selected grid point.
Since $E_s(x)\le e^x$ at nonnegative arguments, $r_t\ge r$.
At the preceding grid point the test fails, while its error in $xe^x$ is
at most
\[
 (n+1)2^{-s}\le2^{-t-8}<\delta/4.
\]
The function $F(x)=xe^x$ has slope at least $1$ on $[0,\infty)$; the
inequality $F(y)-F(x)\ge y-x$ also follows directly from the positive
power series for $F$. Consequently
\[
 0\le r_t-r<\delta+\delta/4<2^{-t}.
\]
All tests here compare finite rational numbers. The algorithm never asks
whether two represented real numbers are equal, or on which side of a
rational a possibly equal real lies. The function \texttt{wApprox} is
therefore a total PR program with the indicated modulus.

\paragraph{Square root and $\pi$.}
The source uses
\[
 R_t(x)=2^{-t}\left\lfloor
 \sqrt{\lfloor x2^{2t}\rfloor}\right\rfloor
\]
for nonnegative $x$. The integer square root may be implemented by bounded
search; its defining square inequalities give
$0\le\sqrt{x}-R_t(x)<2^{-t}$ directly. Lean's \texttt{Nat.sqrt} is used
for this executable presentation, with a bounded-search implementation
available for its PR compilation. The rational approximation to $\pi$ is
the first $2^{t+4}$ terms of
$4\sum_{k\ge0}(-1)^k/(2k+1)$. Pairing adjacent terms proves the alternating
tail estimate $4/(2^{t+5}+1)<2^{-t}$ by finite arithmetic. The usual
identification with $4\arctan(1)$ fixes the normalization of $\pi$.

These choices are deliberately generous. They make the PR bounds explicit
and easy to inspect; they are not intended as efficient numerical
implementations.

\subsection{Exact arithmetic targets for the asymptotic assertions}

The definition \texttt{LogLambertTarget} is the statement that, for all
$n\ge3$ and all $t$, the rational approximation to
\[
 \left|\frac{\log B_n}{n}-
 \left(r-1+\frac1r\right)\right|
\]
is bounded by $(L_t(n)+4)/n+8\varepsilon_t$, using $r_t$ in place of $r$.
This is a first-order arithmetic statement: every approximate quantity is
a fixed PR term and every comparison clears positive denominators. The
extra $8\varepsilon_t$ accounts for the approximation errors. Indeed,
$r_t\ge r\ge1$ gives
$|1/r_t-1/r|\le\varepsilon_t$, and the logarithmic estimates contribute
at most another $2\varepsilon_t/n$.

For the leading equivalent, use the algebraically simpler expression
\[
 C(n)=\left(\frac nr\right)^n
       \frac{\exp(n/r-n-1)}{\sqrt r}.
\]
It is exactly the Wikipedia/Lov\'asz expression, after cancelling the two
square roots involving $n$. The source \texttt{growthApprox n t} evaluates
this expression at internal precision
$s=t+8(n+1)^2+16$.
Here is an explicit error check for that overprecision. For $n\ge3$,
$r\ge1$, $r_s\ge r$, and $a_s=n/r_s\le a=n/r\le n$ satisfy
$|a_s-a|\le n\varepsilon_s$. Therefore
\[
 |a_s^n-a^n|\le n^{n+1}\varepsilon_s.
\]
Both exponential arguments are negative; the Lipschitz constant of the
exponential there is at most $1$. Thus its approximated value differs
from $\exp(a-n-1)$ by at most $(n+1)\varepsilon_s$, and has magnitude at
most $2$. The approximated square root is at least $1$ and differs from
$\sqrt r$ by at most $3\varepsilon_s/2$. Splitting the difference of the
three factors gives
\[
 |C_t(n)-C(n)|
 \le4(n+1)n^n\varepsilon_s
 \le2^{2n^2}\varepsilon_s
 <2^{-t}.
\]
The elementary inequalities used here all hold for $n\ge3$; in particular
$4(n+1)\le2^{n^2}$ and $n^n\le2^{n^2}$.

To control the final quotient, observe that
\[
 \log C(n)=n(r-1+1/r)-1-\tfrac12\log r
 \ge\tfrac12(n-1)\ge1.
\]
We used $r-1+1/r\ge1$, $\log r\le r-1$, and $r\le n$.
Thus $C(n)>2$. Re-evaluate $C$ at precision
$t+4(n+1)^2+8$; its rational approximation is at least $1$, and
$B_n\le n!\le n^n\le2^{n^2}$ gives
\[
 \left|\frac{B_n}{C_{t+4(n+1)^2+8}(n)}-
              \frac{B_n}{C(n)}\right|<2^{-t}.
\]
Accordingly \texttt{WikipediaGrowthTarget} states the rational inequality
with right side $1/q+4\varepsilon_t$, for
\[
 q\ge1,\qquad n\ge2^{64q^2}.
\]
The saddle-point theorem in the human proof supplies $1/q$; the preceding
finite error calculation supplies the rational slack. This is the precise
arithmetic replacement for the asymptotic equivalence. No unrestricted
quantification over real numbers is being added to $\ipra$.

For the two-correction expansion, the additional definitions are the finite
polynomials $U_j$ for $j\le8$, the rational expressions $P,Q$ displayed in
Section~\ref{sec:corrections}, and the normalized expression
\[
 D(n,h)=\frac{B_{n+h}r^{n+h}\sqrt{2\pi av}}{(n+h)!e^{a-1}}.
\]
The target is the finite-precision form of
$|D(n,h)-1-P(r,h)/a-Q(r,h)/a^2|\le2^{400}(H+1)^6/a^3$,
with the stated hypotheses. A signed integer code represents $h$.
All its extra algebra is finite polynomial multiplication and Gaussian
moments through degree $78$. The proof requires the eighth-order phase
Taylor bound and the sixth derivative of the exponential of a degree-five
polynomial, exactly as displayed in that section. These additional targets
are specified in the paper; unlike the two leading-growth targets, they
are not yet packaged as named definitions in \texttt{BellTargets.lean}.

\subsection{A concrete formalization order}

The outstanding work has the following dependency order.
\begin{enumerate}
\item Implement the fixed PR codes for integer/rational arithmetic,
finite arrays and bounded folds. Prove their defining equations and
positive-denominator closure using \texttt{PRA.precZero},
\texttt{PRA.precSucc}, equality congruence and \texttt{PRA.qfInd}.
\item Prove the packed-table invariant and
\texttt{PackedAgreementTarget}. Compile each finite-sum lemma as a schema
for fixed PR summands. Replace any ambient Lean induction over a
universally quantified family by its bounded Boolean invariant.
\item Formalize the bounded restricted-growth encodings and the explicit
bijections in the paper. Derive the Stirling and Bell recurrences, the
finite summation identities, and the count interpretations.
\item Formalize factorial cancellation and the finite moment identities.
For the determinant, use an explicit permutation-sum determinant and the
triangular change of polynomial basis. For congruences, use remainder
arithmetic, bounded orbit decompositions and the finite shift-polynomial
calculation. The propositions \texttt{TouchardTarget},
\texttt{PrimePeriodTarget}, \texttt{LogConvexTarget} and
\texttt{NormalizedLogConcaveTarget} give principal endpoints.
\item Prove the rational tail estimates above. Represent exponential,
logarithm, trigonometric functions and the particular square roots and
Lambert values by these fixed PR approximation programs. Prove their
algebraic and order laws at finite precision before using real notation
as an abbreviation.
\item For each integral in the analytic proof, write the finite rational
mesh sum, its explicit mesh error, and its truncation error. Prove finite
coefficient orthogonality, the local Taylor bounds and the two tail
bounds. Combine the estimates with the explicit thresholds from the
growth section. This yields the arithmetic targets just stated.
\item Only after constructing these \texttt{Prf PRAAx} derivations, apply
\texttt{pra\_iff\_teq} or \texttt{pra\_iff\_native} to transport them to
the selected concise calculus.
\end{enumerate}

The essential remaining workload is explicit: it is the construction of
the arithmetic derivations, especially the finite-precision integral
estimates. The supplied definitions, rates, and human proofs specify what
those derivations must prove. Neither an unrestricted completeness theorem
nor an unproved transfer from classical real analysis is used to declare
the task complete.
